% fig_vl07_dichte.tex — Warum $\Psi$ keine Dichte ist und $|\Psi|^2$ eine:
% dieselbe Wellenfunktion einmal als Real-/Imaginärteil und einmal als
% Wahrscheinlichkeitsdichte, deren Fläche über einem Intervall die
% Ansprechwahrscheinlichkeit des Detektors $F_{[a,b]}$ ist.
%
% Folgt den Konventionen von fig_vl11_kasten.tex.
%
% Geometrie aus der Formel, nicht freihand.  Beide Panels zeigen ein und
% denselben normierten Zustand — ein Gaußsches Wellenpaket, den einfachsten
% Zustand mit $\Psi(x)\ne\delta(x-y)$, den die Vorlesung fordert:
%
%     Psi(x) = (2 pi sigma^2)^(-1/4) exp(-(x-x0)^2/(4 sigma^2)) exp(i k x),
%
% mit sigma = 1, x0 = 0, k = 3.  Daraus folgt alles Gezeichnete:
%
%   Re Psi  = A exp(-x^2/4) cos(3x),   A = (2 pi)^(-1/4) = 0.63162,
%   Im Psi  = A exp(-x^2/4) sin(3x),
%   |Psi|   = A exp(-x^2/4)                      (die Einhüllende),
%   p(x)    = |Psi(x)|^2 = exp(-x^2/2)/sqrt(2 pi),  Maximum 0.39894.
%
% Die Trägerwelle exp(ikx) hebt sich im Betragsquadrat exakt weg: die
% Oszillation in (a) hat in (b) keine Spur hinterlassen.  Genau das ist
% in Prosa mühsam und im Bild eine Sekunde.
%
% Die eingefärbte Fläche ist keine Schätzung, sondern
%     P[x in [a,b]] = Phi(1.5) - Phi(-0.5) = 0.62466,
% mit a = -0.5, b = 1.5 und der Standardnormalverteilung Phi.
% Der schmale Streifen bei c = -1.8 hat Breite eps = 0.2 und damit
% Fläche p(c) eps = 0.0158: dieselbe Konstruktion, nur schmaler.  Er ist
% der Grund, warum F_{x} = |x><x| kein Operator ist — schrumpft eps auf
% null, schrumpft die Fläche mit.
%
% Entwurfsentscheidungen:
%   * Farbe trägt genau zwei Bedeutungen.  In (a) unterscheidet sie die
%     beiden reellen Komponenten desselben komplexen Objekts (Blau/Orange);
%     in (b) ist alles, was Wahrscheinlichkeit ist, blau — der breite und
%     der schmale Streifen unterscheiden sich nur in der Sättigung, weil
%     sie dieselbe Größe sind, nur verschieden breit.
%   * Referenzen sind neutral grau: die Einhüllende +-|Psi| in (a) ist
%     gepunktet grau, keine Daten.
%   * Die beiden Wahrscheinlichkeiten sind gefüllt und nicht schraffiert:
%     schraffiert wird hier, was verboten ist; hier ist die Fläche selbst
%     die Größe, die man ablesen soll, und eine Schraffur würde ihren
%     Inhalt (die Integralformel) unlesbar machen.
%   * Im Sinne des Graustufendrucks ist Im Psi gestrichelt und Re Psi
%     durchgezogen; die beiden Flächen unterscheiden sich zusätzlich
%     durch ihre Breite, nicht nur durch den Grauwert.
%   * Beide Panels benutzen denselben x-Bereich und denselben Maßstab,
%     damit die Breiten vergleichbar sind: die Einhüllende von Psi ist eine
%     Gaußkurve der Breite sqrt(2) sigma, die Dichte p eine der Breite
%     sigma.  Das Quadrieren verschmälert die Glocke um genau diesen
%     Faktor sqrt(2) — im Bild sofort zu sehen, in Prosa eine Fußnote.
%   * Der Streifen sitzt bei c = -1.8 und nicht dichter am Intervall,
%     damit er nicht als dessen Rand missverstanden wird.

\begin{center}
\begin{tikzpicture}

% ─────────────────────────────────────────────────────────────
%  (a)  Die Wellenfunktion selbst: komplex, vorzeichenbehaftet
% ─────────────────────────────────────────────────────────────
\begin{axis}[
    name=psi,
    tpcplot,
    width=6.6cm, height=6.2cm,
    grid=none,
    axis x line=middle,
    axis y line=none,
    axis line style={-{Stealth[length=5pt]}, spacecadet, line width=0.6pt},
    xtick=\empty,
    xmin=-3.6, xmax=3.9,
    ymin=-0.94, ymax=0.94,
    clip=false,
    domain=-3.6:3.75, samples=401,
    title={(a) $\Psi$ \tpcifenglish{is complex}{ist komplex}},
]

  % Einhüllende +-|Psi| = +-A exp(-x^2/4): Referenz, daher grau
  \addplot[birrengray!85, line width=0.7pt, densely dotted]
    { 0.63162*exp(-x^2/4)};
  \addplot[birrengray!85, line width=0.7pt, densely dotted]
    {-0.63162*exp(-x^2/4)};

  % Real- und Imaginärteil derselben Wellenfunktion
  \addplot[wavefn, cgblue]
    {0.63162*exp(-x^2/4)*cos(deg(3*x))};
  \addplot[birrenorange, line width=1.1pt, densely dashed]
    {0.63162*exp(-x^2/4)*sin(deg(3*x))};

  % Direkte Beschriftung statt Legende: jeweils am eigenen Extremum.
  % Re Psi hat sein Maximum A bei x = 0, Im Psi sein Minimum -0.58978
  % bei x = -pi/6 = -0.5236.
  \node[lbl, text=cgblue, anchor=south]      at (axis cs:0,0.66)       {$\real\Psi$};
  \node[lbl, text=birrenorange, anchor=north] at (axis cs:-0.524,-0.68) {$\imag\Psi$};

  % Einhüllende ansprechen, ohne sie zu berühren
  \node[lbl, text=birrengray, font=\scriptsize, anchor=west]
    at (axis cs:-3.55,0.44) {$\pm|\Psi|$};
  \draw[birrengray!85, line width=0.4pt]
    (axis cs:-2.90,0.42) -- (axis cs:-2.45,0.155);

  \node[lbl, text=birrengray, font=\scriptsize, anchor=east]
    at (axis cs:-3.62,0.02) {$0$};

  % Achsenbeschriftung als Knoten unter der Pfeilspitze, nicht als xlabel:
  % ein xlabel ragt rechts aus dem Rahmen und stößt im 12pt-Satz gegen
  % die y-Beschriftung von (b).
  \node[lbl, anchor=north] at (axis cs:3.55,-0.06) {$x$};
\end{axis}

% ─────────────────────────────────────────────────────────────
%  (b)  Das Betragsquadrat: Fläche = Wahrscheinlichkeit
% ─────────────────────────────────────────────────────────────
\begin{axis}[
    tpcplot,
    at={(psi.south east)}, anchor=south west, xshift=22mm,
    width=6.6cm, height=6.2cm,
    grid=none,
    ylabel={$p(x)$},
    every axis y label/.style={at={(axis description cs:-0.17,0.5)},
                               rotate=90, anchor=south, font=\small},
    xmin=-3.6, xmax=3.9,
    ymin=0, ymax=0.52,
    xtick={-1.8,-0.5,1.5}, xticklabels={$c$,$a$,$b$},
    ytick={0,0.1,0.2,0.3,0.4},
    clip=false,
    domain=-3.6:3.75, samples=301,
    title={(b) $p=|\Psi|^{2}$ \tpcifenglish{is the density}{ist die Dichte}},
]

  % P[x in [a,b]] = Flaeche unter p zwischen a und b
  \addplot[draw=none, fill=cgblue!20, domain=-0.5:1.5, samples=101]
    {exp(-x^2/2)/sqrt(2*pi)} \closedcycle;
  \draw[cgblue!55, line width=0.5pt] (axis cs:-0.5,0) -- (axis cs:-0.5,0.35207);
  \draw[cgblue!55, line width=0.5pt] (axis cs: 1.5,0) -- (axis cs: 1.5,0.12952);

  % Derselbe Schnitt, nur um eps = 0.2 breit: gleiche Farbe, dunkler,
  % weil es dieselbe Größe ist.
  \addplot[draw=none, fill=cgblue!60, domain=-1.9:-1.7, samples=11]
    {exp(-x^2/2)/sqrt(2*pi)} \closedcycle;

  % Die Dichte
  \addplot[spacecadet, line width=1.1pt] {exp(-x^2/2)/sqrt(2*pi)};

  % Was die Fläche ist, steht in der Fläche
  \node[lbl, anchor=center] at (axis cs:0.5,0.15) {$\int_a^b p\,dx$};

  % Die Fläche beziffert: die Aussage steht über der Fläche, die sie meint
  \node[lbl, anchor=center] at (axis cs:0.5,0.455)
    {$\mathbb{P}\bigl[x\in[a,b]\bigr]=\tpcifenglish{0.62}{0{,}62}$};

  \node[lbl, anchor=north] at (axis cs:3.55,-0.012) {$x$};

  % Der schmale Streifen und seine Pointe
  \node[lbl, text=spacecadet, font=\scriptsize, anchor=south]
    at (axis cs:-1.8,0.095) {$\varepsilon$};
  \node[lbl, font=\scriptsize, anchor=west, align=left] at (axis cs:-3.40,0.275)
    {$\varepsilon\to0$\\[-1pt]$\Rightarrow\mathbb{P}\to0$};
  \draw[birrengray!85, line width=0.4pt]
    (axis cs:-2.14,0.225) -- (axis cs:-1.90,0.130);
\end{axis}

\end{tikzpicture}
\end{center}
